% ============================================================================
% REVISED RESULTS — paste into manuscript. Real leak-free numbers.
% [[ECE]] etc. = to finalize from the calibration run.
% ============================================================================

\section{Results}
\label{results}
All results use strict subject-level, leakage-free cross-validation (5 folds
$\times$ 10 repeats), with ROC-AUC as the primary metric (the binary task has
$\approx$20\% positives, so accuracy is misleading) and 95\% confidence intervals
across folds. Comparisons across feature sets use identical CV splits (paired).

\subsection{Foundation-model EEG vs.\ hand-crafted EEG}
On MIPDB ($N=118$; 23 positive), frozen LaBraM EEG embeddings substantially
outperform the hand-crafted five-channel features. With mean+std pooling, LaBraM
reaches ROC-AUC $0.660 \pm 0.031$ versus $0.429 \pm 0.037$ for hand-crafted EEG
(paired $\Delta = +0.231 \pm 0.050$; mean-pooling: $0.634$ vs.\ $0.466$,
$\Delta=+0.167$). This reverses the original submission's finding: with a modern
foundation-model representation, EEG moves from the weakest to the strongest
modality, directly answering whether EEG is worth including.

\subsection{Single-modality and multimodal comparison}
Table~\ref{tab:headtohead} reports all modalities and fusion variants under the
identical protocol. Eye-tracking ($0.51$) and behaviour ($0.40$) are at or below
chance in this cohort. Critically, \emph{no} fusion strategy exceeds EEG-alone:
naive concatenation ($0.48$--$0.54$), attention-based set fusion ($0.59$), and
stacking late fusion ($0.47$) all fall at or below the $0.66$ EEG-only ceiling. A
simulated-missingness sweep confirms the model relies almost entirely on EEG
(dropping EEG at test $\rightarrow$ chance; dropping behaviour slightly
\emph{improves} performance). We therefore report the multimodal result honestly:
in MIPDB the non-EEG modalities add no transdiagnostic signal, and we soften the
three-modality claim accordingly.

\begin{table}[t]
\centering\footnotesize
\caption{Leak-free transdiagnostic screening on MIPDB ($N=118$), ROC-AUC
(mean $\pm$ 95\% CI), mean+std pooling.}
\label{tab:headtohead}
\begin{tabular}{lc}
\toprule
\textbf{Feature set / model} & \textbf{ROC-AUC} \\
\midrule
Hand-crafted EEG (Hjorth + spectral entropy) & $0.43 \pm 0.04$ \\
Eye-tracking & $0.51 \pm 0.05$ \\
Behaviour & $0.40 \pm 0.05$ \\
Naive concatenation & $0.48 \pm 0.04$ \\
NeuroFormer attention fusion & $0.59 \pm 0.04$ \\
Stacking late fusion & $0.47 \pm 0.06$ \\
\textbf{Foundation-model EEG (LaBraM)} & $\mathbf{0.66 \pm 0.03}$ \\
\bottomrule
\end{tabular}
\end{table}

\subsection{Calibration and referral}
On the leak-free out-of-fold scores, temperature scaling reduces the Expected
Calibration Error ($[[\text{ECE}_\text{raw}]] \rightarrow [[\text{ECE}_\text{cal}]]$),
and split-conformal prediction at $\alpha{=}0.2$ yields the target coverage with
$[[\,]]\%$ decisive (singleton) predictions; the remainder are referred for
clinician review. The risk--coverage curve shows accuracy increasing as
low-confidence cases are deferred, supporting use as a triage tool.

\subsection{External validation on HBN}
We externally validated the frozen-EEG representation on the independent Healthy
Brain Network cohort at well-powered scale ($N=1{,}980$, RestingState). Predicting
the CBCL psychopathology factors, the general $p$-factor is at chance
(ROC-AUC $0.497 \pm 0.006$; Pearson $r=0.00$); the specific factors are at most
marginally above (internalizing $0.55$, $r=0.09$; attention $0.53$; externalizing
$0.52$), and direct MIPDB$\rightarrow$HBN transfer is at chance ($0.51$). An earlier
underpowered estimate ($N=128$) suggested $\approx0.59$, but this did not survive
scaling and is reported as small-sample variability. Thus the within-MIPDB EEG
signal does \emph{not} generalize to this cohort under a stringent out-of-
distribution test that differs in label construct (clinical diagnosis vs.\ CBCL
factor) and paradigm (mixed vs.\ resting-state). We present this as an honest,
well-powered characterization of the within-cohort--to--OOD gap rather than a
positive generalization claim.

\subsection{Summary}
Foundation-model EEG is the locus of the transdiagnostic signal within MIPDB and
clearly surpasses hand-crafted EEG; additional modalities do not improve screening
here; and the signal does not generalize out-of-distribution at scale. These
leakage-free, externally tested results constitute a rigorous, cautionary
proof-of-concept.
